% SPDX-License-Identifier: GPL-3.0-or-later
% {{NAME}} ({{NATIVE}}) -- what the reading editions' preamble leaves to the language.
% Input by lib/frank-preamble.tex as \input{\FrankLib/lang/\BookLang.tex};
% docs/languages.md section 6 lists the hooks every file here must define and
% section 12 the ones it may; the guide's page "Adding a language"
% (html-guide/markdown/lookup-and-languages/adding-a-language.md) walks
% through writing one.
%
% Scaffolded from lib/lang/_template.tex by lib/newlang.py.  Everything below
% the switches came out of lib/languages.json, so it agrees with the registry
% by construction; \FrankHowTo did not, and is the one thing you must write.
% The model to read while writing it is lib/lang/fa.tex.

% ---- the babel language and the switches ----------------------------------
% \FrankLangName is what \foreignlanguage and \begin{otherlanguage} are given
% everywhere in the core preamble: babel must know this name, or every passage
% comes out in the roman with a "Unknown language" error.
\newcommand{\FrankLangName}{{{BABEL}}}
{{RTL}}% {{RTL_WHY}}
{{HAS_BARE}}% {{HAS_BARE_WHY}}
{{HAS_ALT}}% whether \FrankAltFont exists (the title, the chapter number)
{{HAS_READING}}% whether \chr's third argument is set as a reading
{{HAS_VERT}}% whether pass 4 is \FrankVertPass rather than the font pass
{{HAS_ALOUD}}% whether the chunks carry words, and pass 2 is each one's reading alone

% ---- fonts ------------------------------------------------------------------
% No face is NAMED here: \FrankPickFont walks the registry's tex.main and then
% tex.main_fallbacks and takes the first one installed, so the .tex and the
% .json cannot drift apart (docs/languages.md section 12).  A bundled face --
% one that travels in lib/fonts/ -- is found because build.sh puts that
% directory on OSFONTDIR.
{{FONT_BLOCK}}

% ---- the vocabulary line -----------------------------------------------------
% \vb{{{VB_FORM1}}}{sound}{{{VB_FORM2}}}{sound}{{{VB_FORM3}}}{sound}{meaning}:
% these two are the labels printed before the second and third forms.  The
% three forms are the registry's "vb_forms", which the chunk editor names
% when somebody types a \vb:
%   {{VB_FORMS}}
% Whatever you choose here, docs/lang/{{CODE}}.md must name the same three
% forms in the same order, or the annotator fills slots the page labels
% differently.
% The core prints the second and the third pair only when its FORM is given
% (lib/frank-preamble.tex), so a verb that has no such form leaves the pair
% blank, {}{}, rather than repeating another one.  What the three cannot say
% -- an auxiliary, a class, a governed case -- goes in ONE parenthesis after
% the meaning, items parted by "; ", each only where it is not the ordinary
% case: {to drive (er fährt; aux. sein)} is German's.
% TODO: say here which forms this language gives and why those three, and
% which extras it has, if any; lib/lang/tr.tex and lib/lang/zh.tex are the
% models for a language whose forms are not a Western paradigm.
% The pair below comes from the registry's "vb_labels" ({{VB_LABELS}}), which is
% where the reader takes it from as well: change one and change the other, or
% the PDF and the reader of one book label the same form differently.
% lib/newlang.py --check compares them.
\newcommand{\FrankVbPres}{{{VB_PRES}}}
\newcommand{\FrankVbPast}{{{VB_PAST}}}

% ---- the two Lua filters ---------------------------------------------------
% frank_strip(s): the bare form of a chunk -- what pass 3 shows.  Written from
% the registry's `strip`; the identity when there is nothing to take off.
% frank_latin(s): the label's digits as ASCII, because a PDF destination name
% cannot carry the language's own figures.  Written from `digits`.
% Both are called for every language, so both must exist even when they do
% nothing.  No #, % or literal backslash in a \directlua body: TeX reads it
% first (NOTES section 9 of the Persian book).
\directlua{
{{LUA_FILTERS}}
}

{{TODO}}% ---- the "How to read this" page -----------------------------------------------
% Printed by lib/frank-frontmatter.tex, before the first chapter: the only
% place the reader is told what the passes are for.  Write it in the second
% person, one paragraph per pass, then a paragraph on the transliteration
% scheme and one on how verbs are given -- the shape lib/lang/fa.tex has.
% \pw{...} sets one word of the language inside the English; it is defined by
% the core preamble, after this file, and only expanded when the page is set.
%
% TODO: this is a placeholder.  A passage of {{NAME}} comes {{NPASSES}}:
% {{PASSES}}.
\newcommand{\FrankHowTo}{%
Every passage comes {{NPASSES}}.

\smallskip
\textbf{First} the whole sentence, and nothing else. Try to read it. Get as
far as you can and do not worry about what you cannot get --- the point is
the attempt, made before any help arrives.

\smallskip
\textbf{Second} the same sentence in small chunks, {{NAME}} on one side and on
the other three lines: how it sounds, then how to look each word up, then what
it means. Now you find out how much you had right.

\smallskip
TODO: one paragraph for each remaining pass, then the transliteration scheme
(what the letters with marks stand for) and how verbs are given: as
{{VB_FORM1}}, {{VB_FORM2}}, {{VB_FORM3}}, meaning, and what stands in brackets after it.
}
